%! The Secant, Cosecant, and Cotangent Functions — Unit 3, Lesson 3.11 — Group Activity
\documentclass[10pt]{article}
\usepackage{precalculus-article}
\usepackage{precalculus-boxes}
\newcommand{\TallMath}[1]{$\displaystyle #1\rule[-1.4em]{0pt}{3.2em}$}

\begin{document}

\pageheader{Unit 3, Lesson 3.11}{Group Activity}
\namepartnerperiod

\vspace{0.10in}

%----------------------------------------------------------------------
% Context
%----------------------------------------------------------------------
\begin{scenariobox}[Context: Reciprocal Trigonometric Functions]{sky}
The secant, cosecant, and cotangent functions are built from the reciprocals
of cosine, sine, and tangent. In this activity your group will compute values
of these functions, explore their domains and asymptotes, and investigate
how the graph of a reciprocal function relates to its parent. Choose
\textbf{one tier} based on your readiness level.
\end{scenariobox}

\vspace{0.08in}

%----------------------------------------------------------------------
% Tier R
%----------------------------------------------------------------------
\begin{tcolorbox}[colback=white, colframe=black!40,
    title=\textbf{Tier R --- Reinforce},
    fonttitle=\bfseries, arc=2mm, left=3mm, right=3mm, top=2mm, bottom=2mm]

\textbf{Directions:} Use the given values of $\sin\theta$ and $\cos\theta$ to
compute $\sec\theta$, $\csc\theta$, and $\cot\theta$. Write ``und'' if the
function is undefined at that angle.

\vspace{4pt}
\renewcommand{\arraystretch}{1.7}
\begin{tabular}{|c|c|c|c|c|c|}
\hline
$\theta$ & $\sin\theta$ & $\cos\theta$ & $\sec\theta = \frac{1}{\cos\theta}$ & $\csc\theta = \frac{1}{\sin\theta}$ & $\cot\theta = \frac{\cos\theta}{\sin\theta}$ \\
\hline
$0$ & $0$ & $1$ & & & \\
\hline
$\dfrac{\pi}{6}$ & $\dfrac{1}{2}$ & $\dfrac{\sqrt{3}}{2}$ & & & \\
\hline
$\dfrac{\pi}{4}$ & $\dfrac{\sqrt{2}}{2}$ & $\dfrac{\sqrt{2}}{2}$ & & & \\
\hline
$\dfrac{\pi}{3}$ & $\dfrac{\sqrt{3}}{2}$ & $\dfrac{1}{2}$ & & & \\
\hline
$\dfrac{\pi}{2}$ & $1$ & $0$ & & & \\
\hline
$\pi$ & $0$ & $-1$ & & & \\
\hline
$\dfrac{3\pi}{2}$ & $-1$ & $0$ & & & \\
\hline
$2\pi$ & $0$ & $1$ & & & \\
\hline
\end{tabular}

\vspace{8pt}
\textbf{1.}\quad List all angles in the table where $\sec\theta$ is undefined:

\vspace{4pt}
$\theta = $ \underline{\hspace{5cm}}

\vspace{6pt}
\textbf{2.}\quad List all angles in the table where $\csc\theta$ is undefined:

\vspace{4pt}
$\theta = $ \underline{\hspace{5cm}}

\vspace{6pt}
\textbf{3.}\quad For the angles where $\csc\theta$ is defined, is $|\csc\theta| \geq 1$
always? Circle: \quad \textbf{Yes} \quad / \quad \textbf{No}

\vspace{6pt}
\textbf{4.}\quad At $\theta = \pi/4$, is $\cot\theta$ positive or negative?
\underline{\hspace{2cm}}
At $\theta = 3\pi/4$, is $\cot\theta$ positive or negative?
\underline{\hspace{2cm}}

\end{tcolorbox}

\vspace{0.08in}

%----------------------------------------------------------------------
% Tier A
%----------------------------------------------------------------------
\begin{tcolorbox}[colback=white, colframe=black!40,
    title=\textbf{Tier A --- Apply},
    fonttitle=\bfseries, arc=2mm, left=3mm, right=3mm, top=2mm, bottom=2mm]

\textbf{Directions:} Answer each question without relying on a table. Show
your reasoning using definitions.

\vspace{6pt}
\textbf{1.}\quad State the domain of each function by listing the values excluded
from the domain.

\vspace{4pt}
(a)\quad $\sec\theta$: excluded $\theta$ values (general form):
\underline{\hspace{4.5cm}}

\vspace{4pt}
(b)\quad $\csc\theta$: excluded $\theta$ values (general form):
\underline{\hspace{4.5cm}}

\vspace{4pt}
(c)\quad $\cot\theta$: excluded $\theta$ values (general form):
\underline{\hspace{4.5cm}}

\vspace{8pt}
\textbf{2.}\quad List all vertical asymptotes of each function in the interval $[0, 2\pi]$.

\vspace{4pt}
$\sec\theta$: \underline{\hspace{3.5cm}} \quad
$\csc\theta$: \underline{\hspace{3.5cm}} \quad
$\cot\theta$: \underline{\hspace{3.5cm}}

\vspace{8pt}
\textbf{3.}\quad Evaluate each expression. Show your work.

\vspace{4pt}
(a)\quad $\sec\!\left(\dfrac{2\pi}{3}\right)$
\hspace{0.4cm} [$\cos(2\pi/3) = -1/2$]
\hspace{1cm} $\sec(2\pi/3) = $ \underline{\hspace{2.5cm}}

\vspace{6pt}
(b)\quad $\csc\!\left(\dfrac{5\pi}{6}\right)$
\hspace{0.4cm} [$\sin(5\pi/6) = 1/2$]
\hspace{1cm} $\csc(5\pi/6) = $ \underline{\hspace{2.5cm}}

\vspace{6pt}
(c)\quad $\cot\!\left(\dfrac{5\pi}{4}\right)$
\hspace{0.4cm} [$\sin(5\pi/4) = -\sqrt{2}/2$, $\cos(5\pi/4) = -\sqrt{2}/2$]
\hspace{0.5cm} $\cot(5\pi/4) = $ \underline{\hspace{2.5cm}}

\vspace{8pt}
\textbf{4.}\quad A student says ``$\sec\theta$ could equal 0.8 for some $\theta$.''
Is the student correct? Use the definition to explain why or why not.

\writelines{2}

\end{tcolorbox}

\vspace{0.08in}

%----------------------------------------------------------------------
% Tier E
%----------------------------------------------------------------------
\begin{tcolorbox}[colback=white, colframe=black!40,
    title=\textbf{Tier E --- Extend},
    fonttitle=\bfseries, arc=2mm, left=3mm, right=3mm, top=2mm, bottom=2mm]

\textbf{Directions:} Investigate the relationship between a function and its
reciprocal.

\vspace{6pt}
\textbf{1. Arcs of Secant.}

The graph of $y = \sec\theta$ shows ``U-shaped'' arcs between consecutive
asymptotes. The vertices of each arc occur at angles where $|\cos\theta| = 1$.

\vspace{4pt}
(a)\quad At what angles in $[0, 2\pi]$ does $|\cos\theta| = 1$?

\vspace{4pt}
$\theta = $ \underline{\hspace{4cm}}

\vspace{4pt}
(b)\quad What are the corresponding values of $\sec\theta$ at those angles?

\vspace{4pt}
$\sec\theta = $ \underline{\hspace{4cm}}

\vspace{4pt}
(c)\quad These are the vertex values of the arcs. Explain: why does an arc of
$y = \sec\theta$ open \emph{away} from the $\theta$-axis at each vertex
(rather than crossing it)?

\writelines{2}

\vspace{6pt}
\textbf{2. Algebraic Verification.}

Show algebraically that $\cot\theta = \dfrac{\cos\theta}{\sin\theta}$
starting from the definitions $\tan\theta = \dfrac{\sin\theta}{\cos\theta}$
and $\cot\theta = \dfrac{1}{\tan\theta}$.

\writelines{3}

\vspace{6pt}
\textbf{3. Near an Asymptote.}

As $\theta \to 0^+$, $\sin\theta \to 0^+$. Describe what happens to
$\cot\theta = \cos\theta/\sin\theta$ as $\theta \to 0^+$.
Use the table values from Tier R (or compute your own) to support your answer.

\writelines{3}

\end{tcolorbox}

\end{document}
